\documentclass[11pt,a4paper]{article}
% Preámbulo genérico para informes y documentos (skill informe-latex).
% Uso: copiar este archivo al proyecto como plantilla/preambulo.tex y, desde el documento,
%   \documentclass[11pt,a4paper]{article}
%   % Preámbulo genérico para informes y documentos (skill informe-latex).
% Uso: copiar este archivo al proyecto como plantilla/preambulo.tex y, desde el documento,
%   \documentclass[11pt,a4paper]{article}
%   % Preámbulo genérico para informes y documentos (skill informe-latex).
% Uso: copiar este archivo al proyecto como plantilla/preambulo.tex y, desde el documento,
%   \documentclass[11pt,a4paper]{article}
%   \input{plantilla/preambulo}      (ruta relativa al .tex)
\usepackage[T1]{fontenc}
\usepackage[utf8]{inputenc}
\usepackage[spanish,es-noshorthands,es-tabla]{babel}
\usepackage{lmodern}
\usepackage{microtype}
\usepackage[margin=2.3cm]{geometry}
\usepackage{xcolor}
\usepackage{booktabs}
\usepackage{tabularx}
\usepackage{array}
\usepackage{enumitem}
\usepackage{tcolorbox}
\usepackage{titlesec}
\usepackage{amsmath,amssymb}
\usepackage{graphicx}
\usepackage{caption}
\usepackage{tikz}
\usepackage[hang,flushmargin]{footmisc}
\usepackage[hidelinks]{hyperref}

\usetikzlibrary{babel,arrows.meta,positioning,calc,fit,decorations.pathreplacing}
\tcbuselibrary{skins,breakable}

\definecolor{azul}{HTML}{1F4E79}
\definecolor{gris}{HTML}{6B7280}
\definecolor{ambar}{HTML}{B45309}
\definecolor{fondoazul}{HTML}{F3F7FB}
\definecolor{fondogris}{HTML}{F4F4F5}
\definecolor{fondoambar}{HTML}{FFF7ED}

\titleformat{\section}{\Large\bfseries\color{azul}}{\thesection}{0.6em}{}
\titleformat{\subsection}{\large\bfseries}{\thesubsection}{0.6em}{}
\setcounter{tocdepth}{2}

% --- Opciones de estructura (llamar en el preámbulo del documento, antes de \begin{document}) ---
% Secciones sin número visible (siguen entrando al índice si se usan sin asterisco).
\newcommand{\sinnumeros}{\setcounter{secnumdepth}{0}}
% Cada \section empieza en página nueva.
\newcommand{\sectionbreak}{}
\newcommand{\seccionesenpaginanueva}{\renewcommand{\sectionbreak}{\clearpage}}
% Índice con hipervínculos (después de \maketitle). No fuerza página nueva entre entradas.
\newcommand{\indice}{{\renewcommand{\sectionbreak}{\addvspace{1.5em}}%
  \hypersetup{linktoc=all}\tableofcontents}}

% --- Cajas ---
% Recuadro neutro (gris): consignas, requisitos, datos de entrada, texto citado largo.
%   \begin{recuadro}{Consigna} ... \end{recuadro}
\newtcolorbox{recuadro}[1]{
  breakable, enhanced, colback=fondogris, colframe=gris,
  boxrule=0.4pt, left=6pt, right=6pt, top=4pt, bottom=4pt,
  title={\bfseries #1}, coltitle=white, colbacktitle=gris, fonttitle=\small, fontupper=\small}
% Destacado (azul): conclusión, resumen ejecutivo, definición clave.
\newtcolorbox{destacado}[1]{
  breakable, enhanced, colback=fondoazul, colframe=azul,
  boxrule=0.5pt, left=6pt, right=6pt, top=5pt, bottom=5pt,
  title={\bfseries #1}, coltitle=white, colbacktitle=azul, fonttitle=\small}
% Advertencias, supuestos, ambigüedades, erratas (ámbar).
\newtcolorbox{aviso}[1]{
  breakable, enhanced, colback=fondoambar, colframe=ambar,
  boxrule=0.5pt, left=6pt, right=6pt, top=5pt, bottom=5pt,
  title={\bfseries #1}, coltitle=white, colbacktitle=ambar, fonttitle=\small}
% Ampliación opcional (franja azul). No es una tcolorbox: las notas al pie dentro de
% una tcolorbox que se parte entre páginas se superponen con el contenido.
%   \begin{extras} ... \end{extras}      o      \begin{extras}[PARA PROFUNDIZAR] ... \end{extras}
\newenvironment{extras}[1][EXTRAS]{%
  \par\medskip\noindent
  {\setlength{\fboxsep}{3pt}\colorbox{azul}{\color{white}\bfseries\small\ #1\ }}%
  \par\nobreak\vspace{1pt}{\color{azul}\hrule height 0.6pt}\nobreak\vspace{3pt}}
 {\par\vspace{3pt}{\color{azul}\hrule height 0.6pt}\medskip}

\setlist[itemize]{itemsep=2pt, topsep=3pt, parsep=0pt}
\setlist[enumerate]{itemsep=2pt, topsep=3pt, parsep=0pt}
\captionsetup{font=small, labelfont=bf}
\setlength{\parindent}{0pt}
\setlength{\parskip}{4pt}

% --- Citas y fuentes ---
% Cita textual en el idioma original, sin traducir.
\newcommand{\cita}[1]{``\textit{#1}''}
% Fuente en nota al pie (nunca en el cuerpo):  \fuente{Autor, \textit{Título}, §2.1, p.~34}
% Varias fuentes en una nota: separarlas con «;».
\newcommand{\fuente}[1]{\footnote{#1.}}
% Nota inicial que explica abreviaturas de fuentes:  \notafuentes{\textbf{AIMA}: Russell \& Norvig, ...}
\newcommand{\notafuentes}[1]{{\small\color{gris} Las fuentes van en notas al pie. #1\par}}

% Íconos reutilizables (pics de TikZ). Uso: \pic[scale=0.8] at (x,y) {robot};
\tikzset{
  pics/robot/.style={code={
    \draw[thick, fill=azul!15] (-0.35,0.15) rectangle (0.35,0.65);            % cabeza
    \fill[azul] (-0.15,0.42) circle (0.06) (0.15,0.42) circle (0.06);         % ojos
    \draw[thick] (0,0.65) -- (0,0.8); \fill[ambar] (0,0.83) circle (0.05);    % antena
    \draw[thick, fill=azul!15] (-0.45,-0.45) rectangle (0.45,0.12);           % cuerpo
    \draw[thick] (-0.45,-0.1) -- (-0.65,0.1) (0.45,-0.1) -- (0.65,0.1);       % brazos
    \draw[thick, fill=gris!40] (-0.25,-0.55) circle (0.12) (0.25,-0.55) circle (0.12); % ruedas
  }},
  pics/compu/.style={code={
    \draw[thick, fill=azul!15, rounded corners=1pt] (-0.55,-0.25) rectangle (0.55,0.5); % pantalla
    \node[font=\tiny\ttfamily\bfseries, text=azul] at (0,0.12) {</>};
    \draw[thick] (0,-0.25) -- (0,-0.4); \draw[thick] (-0.25,-0.4) -- (0.25,-0.4);       % pie
  }},
  pics/persona/.style={code={
    \draw[thick] (0,0.35) circle (0.1);
    \draw[thick] (0,0.25) -- (0,-0.05) (-0.15,0.15) -- (0.15,0.15)
                 (0,-0.05) -- (-0.12,-0.3) (0,-0.05) -- (0.12,-0.3);
  }},
  pics/arbol/.style={code={
    \draw[thick, brown!60!black] (0,-0.3) -- (0,0);
    \fill[green!50!black!70] (0,0.15) circle (0.2);
  }},
  pics/documento/.style={code={
    \draw[thick, fill=white] (-0.2,-0.28) -- (-0.2,0.28) -- (0.1,0.28) -- (0.2,0.18) -- (0.2,-0.28) -- cycle;
    \foreach \y in {0.1,0,-0.1} \draw[gris] (-0.12,\y) -- (0.12,\y);
  }},
  pics/sobre/.style={code={
    \draw[thick, fill=white] (-0.25,-0.17) rectangle (0.25,0.17);
    \draw[thick] (-0.25,0.17) -- (0,-0.02) -- (0.25,0.17);
  }},
}

% --- Macros propios del proyecto ---
\usepackage{needspace}
% Fuentes frecuentes
\newcommand{\nb}[1]{\footnote{NB, §#1.}}
\newcommand{\teo}[1]{\footnote{Teoría, clase #1.}}
      (ruta relativa al .tex)
\usepackage[T1]{fontenc}
\usepackage[utf8]{inputenc}
\usepackage[spanish,es-noshorthands,es-tabla]{babel}
\usepackage{lmodern}
\usepackage{microtype}
\usepackage[margin=2.3cm]{geometry}
\usepackage{xcolor}
\usepackage{booktabs}
\usepackage{tabularx}
\usepackage{array}
\usepackage{enumitem}
\usepackage{tcolorbox}
\usepackage{titlesec}
\usepackage{amsmath,amssymb}
\usepackage{graphicx}
\usepackage{caption}
\usepackage{tikz}
\usepackage[hang,flushmargin]{footmisc}
\usepackage[hidelinks]{hyperref}

\usetikzlibrary{babel,arrows.meta,positioning,calc,fit,decorations.pathreplacing}
\tcbuselibrary{skins,breakable}

\definecolor{azul}{HTML}{1F4E79}
\definecolor{gris}{HTML}{6B7280}
\definecolor{ambar}{HTML}{B45309}
\definecolor{fondoazul}{HTML}{F3F7FB}
\definecolor{fondogris}{HTML}{F4F4F5}
\definecolor{fondoambar}{HTML}{FFF7ED}

\titleformat{\section}{\Large\bfseries\color{azul}}{\thesection}{0.6em}{}
\titleformat{\subsection}{\large\bfseries}{\thesubsection}{0.6em}{}
\setcounter{tocdepth}{2}

% --- Opciones de estructura (llamar en el preámbulo del documento, antes de \begin{document}) ---
% Secciones sin número visible (siguen entrando al índice si se usan sin asterisco).
\newcommand{\sinnumeros}{\setcounter{secnumdepth}{0}}
% Cada \section empieza en página nueva.
\newcommand{\sectionbreak}{}
\newcommand{\seccionesenpaginanueva}{\renewcommand{\sectionbreak}{\clearpage}}
% Índice con hipervínculos (después de \maketitle). No fuerza página nueva entre entradas.
\newcommand{\indice}{{\renewcommand{\sectionbreak}{\addvspace{1.5em}}%
  \hypersetup{linktoc=all}\tableofcontents}}

% --- Cajas ---
% Recuadro neutro (gris): consignas, requisitos, datos de entrada, texto citado largo.
%   \begin{recuadro}{Consigna} ... \end{recuadro}
\newtcolorbox{recuadro}[1]{
  breakable, enhanced, colback=fondogris, colframe=gris,
  boxrule=0.4pt, left=6pt, right=6pt, top=4pt, bottom=4pt,
  title={\bfseries #1}, coltitle=white, colbacktitle=gris, fonttitle=\small, fontupper=\small}
% Destacado (azul): conclusión, resumen ejecutivo, definición clave.
\newtcolorbox{destacado}[1]{
  breakable, enhanced, colback=fondoazul, colframe=azul,
  boxrule=0.5pt, left=6pt, right=6pt, top=5pt, bottom=5pt,
  title={\bfseries #1}, coltitle=white, colbacktitle=azul, fonttitle=\small}
% Advertencias, supuestos, ambigüedades, erratas (ámbar).
\newtcolorbox{aviso}[1]{
  breakable, enhanced, colback=fondoambar, colframe=ambar,
  boxrule=0.5pt, left=6pt, right=6pt, top=5pt, bottom=5pt,
  title={\bfseries #1}, coltitle=white, colbacktitle=ambar, fonttitle=\small}
% Ampliación opcional (franja azul). No es una tcolorbox: las notas al pie dentro de
% una tcolorbox que se parte entre páginas se superponen con el contenido.
%   \begin{extras} ... \end{extras}      o      \begin{extras}[PARA PROFUNDIZAR] ... \end{extras}
\newenvironment{extras}[1][EXTRAS]{%
  \par\medskip\noindent
  {\setlength{\fboxsep}{3pt}\colorbox{azul}{\color{white}\bfseries\small\ #1\ }}%
  \par\nobreak\vspace{1pt}{\color{azul}\hrule height 0.6pt}\nobreak\vspace{3pt}}
 {\par\vspace{3pt}{\color{azul}\hrule height 0.6pt}\medskip}

\setlist[itemize]{itemsep=2pt, topsep=3pt, parsep=0pt}
\setlist[enumerate]{itemsep=2pt, topsep=3pt, parsep=0pt}
\captionsetup{font=small, labelfont=bf}
\setlength{\parindent}{0pt}
\setlength{\parskip}{4pt}

% --- Citas y fuentes ---
% Cita textual en el idioma original, sin traducir.
\newcommand{\cita}[1]{``\textit{#1}''}
% Fuente en nota al pie (nunca en el cuerpo):  \fuente{Autor, \textit{Título}, §2.1, p.~34}
% Varias fuentes en una nota: separarlas con «;».
\newcommand{\fuente}[1]{\footnote{#1.}}
% Nota inicial que explica abreviaturas de fuentes:  \notafuentes{\textbf{AIMA}: Russell \& Norvig, ...}
\newcommand{\notafuentes}[1]{{\small\color{gris} Las fuentes van en notas al pie. #1\par}}

% Íconos reutilizables (pics de TikZ). Uso: \pic[scale=0.8] at (x,y) {robot};
\tikzset{
  pics/robot/.style={code={
    \draw[thick, fill=azul!15] (-0.35,0.15) rectangle (0.35,0.65);            % cabeza
    \fill[azul] (-0.15,0.42) circle (0.06) (0.15,0.42) circle (0.06);         % ojos
    \draw[thick] (0,0.65) -- (0,0.8); \fill[ambar] (0,0.83) circle (0.05);    % antena
    \draw[thick, fill=azul!15] (-0.45,-0.45) rectangle (0.45,0.12);           % cuerpo
    \draw[thick] (-0.45,-0.1) -- (-0.65,0.1) (0.45,-0.1) -- (0.65,0.1);       % brazos
    \draw[thick, fill=gris!40] (-0.25,-0.55) circle (0.12) (0.25,-0.55) circle (0.12); % ruedas
  }},
  pics/compu/.style={code={
    \draw[thick, fill=azul!15, rounded corners=1pt] (-0.55,-0.25) rectangle (0.55,0.5); % pantalla
    \node[font=\tiny\ttfamily\bfseries, text=azul] at (0,0.12) {</>};
    \draw[thick] (0,-0.25) -- (0,-0.4); \draw[thick] (-0.25,-0.4) -- (0.25,-0.4);       % pie
  }},
  pics/persona/.style={code={
    \draw[thick] (0,0.35) circle (0.1);
    \draw[thick] (0,0.25) -- (0,-0.05) (-0.15,0.15) -- (0.15,0.15)
                 (0,-0.05) -- (-0.12,-0.3) (0,-0.05) -- (0.12,-0.3);
  }},
  pics/arbol/.style={code={
    \draw[thick, brown!60!black] (0,-0.3) -- (0,0);
    \fill[green!50!black!70] (0,0.15) circle (0.2);
  }},
  pics/documento/.style={code={
    \draw[thick, fill=white] (-0.2,-0.28) -- (-0.2,0.28) -- (0.1,0.28) -- (0.2,0.18) -- (0.2,-0.28) -- cycle;
    \foreach \y in {0.1,0,-0.1} \draw[gris] (-0.12,\y) -- (0.12,\y);
  }},
  pics/sobre/.style={code={
    \draw[thick, fill=white] (-0.25,-0.17) rectangle (0.25,0.17);
    \draw[thick] (-0.25,0.17) -- (0,-0.02) -- (0.25,0.17);
  }},
}

% --- Macros propios del proyecto ---
\usepackage{needspace}
% Fuentes frecuentes
\newcommand{\nb}[1]{\footnote{NB, §#1.}}
\newcommand{\teo}[1]{\footnote{Teoría, clase #1.}}
      (ruta relativa al .tex)
\usepackage[T1]{fontenc}
\usepackage[utf8]{inputenc}
\usepackage[spanish,es-noshorthands,es-tabla]{babel}
\usepackage{lmodern}
\usepackage{microtype}
\usepackage[margin=2.3cm]{geometry}
\usepackage{xcolor}
\usepackage{booktabs}
\usepackage{tabularx}
\usepackage{array}
\usepackage{enumitem}
\usepackage{tcolorbox}
\usepackage{titlesec}
\usepackage{amsmath,amssymb}
\usepackage{graphicx}
\usepackage{caption}
\usepackage{tikz}
\usepackage[hang,flushmargin]{footmisc}
\usepackage[hidelinks]{hyperref}

\usetikzlibrary{babel,arrows.meta,positioning,calc,fit,decorations.pathreplacing}
\tcbuselibrary{skins,breakable}

\definecolor{azul}{HTML}{1F4E79}
\definecolor{gris}{HTML}{6B7280}
\definecolor{ambar}{HTML}{B45309}
\definecolor{fondoazul}{HTML}{F3F7FB}
\definecolor{fondogris}{HTML}{F4F4F5}
\definecolor{fondoambar}{HTML}{FFF7ED}

\titleformat{\section}{\Large\bfseries\color{azul}}{\thesection}{0.6em}{}
\titleformat{\subsection}{\large\bfseries}{\thesubsection}{0.6em}{}
\setcounter{tocdepth}{2}

% --- Opciones de estructura (llamar en el preámbulo del documento, antes de \begin{document}) ---
% Secciones sin número visible (siguen entrando al índice si se usan sin asterisco).
\newcommand{\sinnumeros}{\setcounter{secnumdepth}{0}}
% Cada \section empieza en página nueva.
\newcommand{\sectionbreak}{}
\newcommand{\seccionesenpaginanueva}{\renewcommand{\sectionbreak}{\clearpage}}
% Índice con hipervínculos (después de \maketitle). No fuerza página nueva entre entradas.
\newcommand{\indice}{{\renewcommand{\sectionbreak}{\addvspace{1.5em}}%
  \hypersetup{linktoc=all}\tableofcontents}}

% --- Cajas ---
% Recuadro neutro (gris): consignas, requisitos, datos de entrada, texto citado largo.
%   \begin{recuadro}{Consigna} ... \end{recuadro}
\newtcolorbox{recuadro}[1]{
  breakable, enhanced, colback=fondogris, colframe=gris,
  boxrule=0.4pt, left=6pt, right=6pt, top=4pt, bottom=4pt,
  title={\bfseries #1}, coltitle=white, colbacktitle=gris, fonttitle=\small, fontupper=\small}
% Destacado (azul): conclusión, resumen ejecutivo, definición clave.
\newtcolorbox{destacado}[1]{
  breakable, enhanced, colback=fondoazul, colframe=azul,
  boxrule=0.5pt, left=6pt, right=6pt, top=5pt, bottom=5pt,
  title={\bfseries #1}, coltitle=white, colbacktitle=azul, fonttitle=\small}
% Advertencias, supuestos, ambigüedades, erratas (ámbar).
\newtcolorbox{aviso}[1]{
  breakable, enhanced, colback=fondoambar, colframe=ambar,
  boxrule=0.5pt, left=6pt, right=6pt, top=5pt, bottom=5pt,
  title={\bfseries #1}, coltitle=white, colbacktitle=ambar, fonttitle=\small}
% Ampliación opcional (franja azul). No es una tcolorbox: las notas al pie dentro de
% una tcolorbox que se parte entre páginas se superponen con el contenido.
%   \begin{extras} ... \end{extras}      o      \begin{extras}[PARA PROFUNDIZAR] ... \end{extras}
\newenvironment{extras}[1][EXTRAS]{%
  \par\medskip\noindent
  {\setlength{\fboxsep}{3pt}\colorbox{azul}{\color{white}\bfseries\small\ #1\ }}%
  \par\nobreak\vspace{1pt}{\color{azul}\hrule height 0.6pt}\nobreak\vspace{3pt}}
 {\par\vspace{3pt}{\color{azul}\hrule height 0.6pt}\medskip}

\setlist[itemize]{itemsep=2pt, topsep=3pt, parsep=0pt}
\setlist[enumerate]{itemsep=2pt, topsep=3pt, parsep=0pt}
\captionsetup{font=small, labelfont=bf}
\setlength{\parindent}{0pt}
\setlength{\parskip}{4pt}

% --- Citas y fuentes ---
% Cita textual en el idioma original, sin traducir.
\newcommand{\cita}[1]{``\textit{#1}''}
% Fuente en nota al pie (nunca en el cuerpo):  \fuente{Autor, \textit{Título}, §2.1, p.~34}
% Varias fuentes en una nota: separarlas con «;».
\newcommand{\fuente}[1]{\footnote{#1.}}
% Nota inicial que explica abreviaturas de fuentes:  \notafuentes{\textbf{AIMA}: Russell \& Norvig, ...}
\newcommand{\notafuentes}[1]{{\small\color{gris} Las fuentes van en notas al pie. #1\par}}

% Íconos reutilizables (pics de TikZ). Uso: \pic[scale=0.8] at (x,y) {robot};
\tikzset{
  pics/robot/.style={code={
    \draw[thick, fill=azul!15] (-0.35,0.15) rectangle (0.35,0.65);            % cabeza
    \fill[azul] (-0.15,0.42) circle (0.06) (0.15,0.42) circle (0.06);         % ojos
    \draw[thick] (0,0.65) -- (0,0.8); \fill[ambar] (0,0.83) circle (0.05);    % antena
    \draw[thick, fill=azul!15] (-0.45,-0.45) rectangle (0.45,0.12);           % cuerpo
    \draw[thick] (-0.45,-0.1) -- (-0.65,0.1) (0.45,-0.1) -- (0.65,0.1);       % brazos
    \draw[thick, fill=gris!40] (-0.25,-0.55) circle (0.12) (0.25,-0.55) circle (0.12); % ruedas
  }},
  pics/compu/.style={code={
    \draw[thick, fill=azul!15, rounded corners=1pt] (-0.55,-0.25) rectangle (0.55,0.5); % pantalla
    \node[font=\tiny\ttfamily\bfseries, text=azul] at (0,0.12) {</>};
    \draw[thick] (0,-0.25) -- (0,-0.4); \draw[thick] (-0.25,-0.4) -- (0.25,-0.4);       % pie
  }},
  pics/persona/.style={code={
    \draw[thick] (0,0.35) circle (0.1);
    \draw[thick] (0,0.25) -- (0,-0.05) (-0.15,0.15) -- (0.15,0.15)
                 (0,-0.05) -- (-0.12,-0.3) (0,-0.05) -- (0.12,-0.3);
  }},
  pics/arbol/.style={code={
    \draw[thick, brown!60!black] (0,-0.3) -- (0,0);
    \fill[green!50!black!70] (0,0.15) circle (0.2);
  }},
  pics/documento/.style={code={
    \draw[thick, fill=white] (-0.2,-0.28) -- (-0.2,0.28) -- (0.1,0.28) -- (0.2,0.18) -- (0.2,-0.28) -- cycle;
    \foreach \y in {0.1,0,-0.1} \draw[gris] (-0.12,\y) -- (0.12,\y);
  }},
  pics/sobre/.style={code={
    \draw[thick, fill=white] (-0.25,-0.17) rectangle (0.25,0.17);
    \draw[thick] (-0.25,0.17) -- (0,-0.02) -- (0.25,0.17);
  }},
}

% --- Macros propios del proyecto ---
\usepackage{needspace}
% Fuentes frecuentes
\newcommand{\nb}[1]{\footnote{NB, §#1.}}
\newcommand{\teo}[1]{\footnote{Teoría, clase #1.}}


\newcommand{\pres}{../../entregables/presentacion-Bertotto-Cid.pdf}
% \slide{n}{título}{tiempo}: encabezado de cada slide con su miniatura
\newcommand{\slide}[3]{%
  \needspace{24\baselineskip}
  \subsection{Slide #1 · #2}
  \begin{center}
    \fcolorbox{gris!50}{white}{\includegraphics[page=#1,width=0.52\linewidth]{\pres}}\\[2pt]
    {\small\color{gris} Tiempo sugerido: \textbf{#3}}
  \end{center}}
\newenvironment{speech}{\begin{destacado}{Qué decir}}{\end{destacado}}

\title{\vspace{-1.5cm}\textbf{Guion de la presentación}\\[3pt]
\large Clasificación de películas con un árbol de decisión y una red neuronal}
\author{\normalsize Tomás Bertotto y Nicolás Cid --- Inteligencia Artificial, UNS --- 2026}
\date{}

\begin{document}
\maketitle
\vspace{-1cm}
\notafuentes{\textbf{NB}: notebook entregado (\texttt{proyecto-Bertotto-Cid.ipynb}), con su número de sección.
\textbf{Teoría}: diapositivas de la cátedra (clases 9, 10 y 11 de Aprendizaje Automático).}

\indice

\section{Cómo usar este guion}

\begin{itemize}
  \item Cada slide tiene: su miniatura, el tiempo sugerido, el texto para decir (caja azul) y, debajo,
        \textbf{Para entenderlo}: conceptos y datos que no se dicen pero conviene saber para defenderlo.
  \item El texto está escrito para decirse en voz alta; no hace falta leerlo literal. Lo importante son
        las ideas en \textbf{negrita}.
  \item Al final hay un banco de preguntas posibles con respuestas cortas, y los puntos débiles del
        trabajo con cómo responderlos.
\end{itemize}

\begin{center}
\small
\begin{tabularx}{\linewidth}{@{}r X r r@{}}
\toprule
\# & Slide & Tiempo & Acumulado \\
\midrule
1  & Portada                               & 15 s & 0:15 \\
2  & Los datos                             & 40 s & 0:55 \\
3  & Problemas encontrados                 & 45 s & 1:40 \\
4  & Colas largas y el logaritmo           & 35 s & 2:15 \\
5  & Qué atributos separan las clases      & 40 s & 2:55 \\
6  & Decisiones y por qué                  & 55 s & 3:50 \\
7  & Sin fuga de datos                     & 40 s & 4:30 \\
8  & El árbol por defecto memoriza         & 40 s & 5:10 \\
9  & Árbol: evolución de los experimentos  & 55 s & 6:05 \\
10 & La popularidad decide primero         & 35 s & 6:40 \\
11 & Cuándo dejar de entrenar              & 45 s & 7:25 \\
12 & Red: evolución de los experimentos    & 55 s & 8:20 \\
13 & Árbol contra red, en validación       & 40 s & 9:00 \\
14 & Elegimos el árbol                     & 35 s & 9:35 \\
\bottomrule
\end{tabularx}
\end{center}

Quedan unos 25 segundos de margen. Si el tiempo aprieta, los lugares para recortar son la slide 5 (dejar solo
la popularidad) y la 12 (no detallar R6 y R7).

\begin{destacado}{La historia en tres frases}
\begin{enumerate}
  \item Los datos tienen trampas (ceros que significan ``no se sabe'', faltantes, colas largas, nombres
        con miles de valores) y las tratamos antes de entrenar, \textbf{sin fuga de datos}.
  \item Los dos modelos \textbf{sobreajustan} si no se los limita; los experimentos consisten en
        controlar eso (profundidad y mínimo para dividir en el árbol; capacidad y épocas en la red).
  \item El \textbf{árbol gana} (0.670 contra 0.644 de F1 macro), sobre todo en \textit{Muy buena}, y
        además es más simple y se puede leer.
\end{enumerate}
\end{destacado}

\clearpage
\section{Guion por slide}

% ------------------------------------------------------------------
\slide{1}{Portada}{15 s}

\begin{speech}
Buenas. Somos Tomás Bertotto y Nicolás Cid. Nuestro proyecto consiste en \textbf{clasificar películas en
tres categorías}: Mala, Buena y Muy buena, a partir de sus datos. Para eso construimos y comparamos
\textbf{dos modelos}: un árbol de decisión y una red neuronal.
\end{speech}

\begin{extras}[PARA ENTENDERLO]
\begin{itemize}
  \item Es \textbf{aprendizaje supervisado}: cada ejemplo de entrenamiento trae su clase, y el modelo
        aprende a predecirla a partir de los atributos.
  \item Es \textbf{clasificación ternaria} (tres clases), no regresión: la salida es una categoría, no un número.
\end{itemize}
\end{extras}

% ------------------------------------------------------------------
\slide{2}{Los datos}{40 s}

\begin{speech}
Partimos de \textbf{12.500 películas etiquetadas} para entrenar y \textbf{500 de test}, para las que
tenemos que predecir la clase. Cada película tiene atributos propios (género, presupuesto, ingresos,
duración, popularidad, productora), datos de sus tres intérpretes principales y de su director.

Lo primero que miramos fue la distribución de la clase: Buena es casi la mitad, Mala un 38\,\%, y
\textbf{Muy buena apenas un 15\,\%}. Este desbalance condiciona cómo evaluamos: con la accuracy sola, un
modelo que casi nunca dice ``Muy buena'' igual puede sacar un buen número. Por eso la métrica principal es
el \textbf{F1 macro}, que promedia el F1 de cada clase y les da a las tres el mismo peso.
\end{speech}

\begin{extras}[PARA ENTENDERLO]
\begin{itemize}
  \item \textbf{Precisión} de una clase: de las veces que el modelo dijo esa clase, cuántas acertó.
        \textbf{Recall}: de los ejemplos que realmente son de esa clase, cuántos encontró.
  \item \textbf{F1} $= \dfrac{2PR}{P+R}$, la media armónica: es baja si cualquiera de las dos es baja.
        Ejemplo real: el árbol en \textit{Muy buena} tiene precisión 0.683 y recall 0.393, y F1 0.499.\nb{4.7}
  \item \textbf{Macro}: promedio simple de los F1 de las tres clases (fila \texttt{macro avg} de
        \texttt{classification\_report}). El \texttt{weighted avg}, en cambio, pondera por cantidad de
        ejemplos, y ahí \textit{Muy buena} pesa poco.
  \item Por qué la accuracy engaña: un modelo que nunca dice \textit{Muy buena} puede acertar hasta el
        85\,\% como máximo, pero su F1 en esa clase es 0 y el F1 macro cae mucho.
\end{itemize}
\end{extras}

% ------------------------------------------------------------------
\slide{3}{Problemas encontrados}{45 s}

\begin{speech}
En el análisis exploratorio encontramos cuatro problemas. El \textbf{79\,\%} de las películas tiene
presupuesto igual a cero, y el \textbf{76\,\%} ingresos igual a cero. Una película no puede costar cero:
\textbf{ese cero significa ``no se sabe''}. Si lo dejáramos como número, el modelo lo trataría como una
película muy barata.

El país de la productora falta en el \textbf{90\,\%} de los casos. Y los nombres de intérpretes,
directores y productoras tienen \textbf{miles de valores distintos}, casi uno por película. Además la
duración falta en un 10\,\%, y las variables de dinero y popularidad tienen colas muy largas, que se ven
en la próxima slide.
\end{speech}

\begin{extras}[PARA ENTENDERLO]
\begin{itemize}
  \item Un \textbf{cero disfrazado} es peor que un faltante explícito: \texttt{isna()} no lo detecta y los
        modelos lo usan como un valor real. Por eso lo buscamos aparte, con \texttt{(train[col] == 0).mean()}.\nb{2.4}
  \item Lo mismo pasa con \texttt{runtime}: una duración de 0 minutos tampoco es posible.
  \item \textbf{Alta cardinalidad}: una variable categórica con muchísimos valores distintos.
\end{itemize}
\end{extras}

% ------------------------------------------------------------------
\slide{4}{Colas largas y el logaritmo}{35 s}

\begin{speech}
Arriba están los histogramas originales de presupuesto, ingresos, popularidad y duración. Casi todo se
amontona en una barra pegada al cero, porque \textbf{unos pocos valores enormes estiran el eje}. Abajo
aplicamos \textbf{logaritmo de uno más x}, sin los ceros, y las distribuciones quedan mucho más simétricas
y comparables. Al árbol esto no le cambia nada, pero a la red neuronal sí, como vamos a ver en el
preprocesamiento.
\end{speech}

\begin{extras}[PARA ENTENDERLO]
\begin{itemize}
  \item Se usa $\log(1+x)$ y no $\log(x)$ porque $\log(0)$ no está definido; con el $+1$, el 0 va a 0.
  \item El logaritmo \textbf{comprime la cola}: pasar de 1 millón a 100 millones es multiplicar por 100,
        pero en escala logarítmica es solo sumar $\log(100) \approx 4{,}6$.
  \item Es \textbf{monótono} (mantiene el orden): por eso al árbol no le cambia nada, porque solo compara
        contra umbrales.
  \item El año no lleva logaritmo: no tiene cola larga.\nb{5.1}
\end{itemize}
\end{extras}

% ------------------------------------------------------------------
\slide{5}{Qué atributos separan las clases}{40 s}

\begin{speech}
También miramos qué atributos se relacionan con la clase. \textbf{La popularidad de la película es el más
claro}: el promedio es 4.9 en las Malas, contra 15.4 en las Buenas y 23.4 en las Muy buenas; en el
boxplot, las Malas quedan claramente abajo. La popularidad del intérprete principal muestra diferencias
chicas, y el año, que sacamos del título, casi no cambia entre clases. Igual lo dejamos: solo no separa,
pero el árbol lo puede combinar con otros atributos. El género también ayuda: por ejemplo, el 69\,\% de
las películas de terror son Malas.
\end{speech}

\begin{extras}[PARA ENTENDERLO]
\begin{itemize}
  \item \textbf{Cómo se lee un boxplot}: la caja va del cuartil 1 al 3 (el 50\,\% central de los datos),
        la línea interna es la mediana, los bigotes llegan hasta 1,5 veces el ancho de la caja y los
        puntos sueltos son valores atípicos.
  \item Los boxplots de popularidad están en escala $\log(1+x)$; si no, la cola larga aplastaría las cajas.
  \item Otro dato del género: \textit{Documentary} es el que tiene más proporción de \textit{Muy buena} (35\,\%).\nb{2.9}
  \item Que esto lo hayamos visto en el análisis exploratorio y después el árbol lo confirme (slide 10)
        es un buen argumento de coherencia del trabajo.
\end{itemize}
\end{extras}

% ------------------------------------------------------------------
\slide{6}{Decisiones y por qué}{55 s}

\begin{speech}
Con esto tomamos estas decisiones. Extrajimos el \textbf{año del título}, porque no venía como columna.
Los ceros de presupuesto, ingresos y duración pasan a ser \textbf{faltantes}, y agregamos un
\textbf{indicador} de si el dato se conocía, porque que la película tenga el presupuesto publicado ya
puede ser información.

Al país no lo completamos: con 90\,\% de faltantes \textbf{estaríamos inventando casi todos los valores};
nos quedamos solo con si se conoce o no. Descartamos identificadores, título y nombres: con miles de
valores, un one-hot generaría miles de columnas casi vacías y el modelo memorizaría nombres; la
\textbf{popularidad numérica de cada persona ya resume esa información}.

El género, que tiene 8 valores sin orden, va con \textbf{one-hot}. Los faltantes numéricos se completan
con el \textbf{promedio}, como hicimos con la edad en el laboratorio 1. Y \textbf{solo para la red}
aplicamos el logaritmo y escalamos al intervalo de menos uno a uno, porque el descenso por gradiente
necesita escalas comparables.
\end{speech}

\begin{extras}[PARA ENTENDERLO]
\begin{itemize}
  \item \textbf{One-hot}: una columna 0/1 por cada valor del género. Se usa en vez de numerar los géneros
        (1, 2, 3\dots) porque eso inventaría un orden y distancias que no existen (``Drama'' no es el
        doble de ``Comedia''). Los faltantes de género pasan a ser la categoría ``Desconocido''.
  \item Resultado: \textbf{28 atributos} (15 numéricos, 4 binarios y las columnas del género).\nb{4.1}
  \item El orden importa: los ceros pasan a faltante \textbf{antes} de calcular el promedio; si no, el
        promedio del presupuesto saldría tirado hacia cero.
  \item \texttt{belong\_to\_collection} (si es parte de una saga) se pasa a 0/1.
  \item \texttt{MinMaxScaler(-1, 1)}: $x' = -1 + 2\,\dfrac{x - \min}{\max - \min}$, con mínimo y máximo
        de cada columna calculados en el entrenamiento.
\end{itemize}
\end{extras}

% ------------------------------------------------------------------
\slide{7}{Sin fuga de datos}{40 s}

\begin{speech}
Un punto que cuidamos especialmente es la \textbf{fuga de datos}. Primero \textbf{dividimos}: 80\,\%
para entrenar y 20\,\% para validar, de forma \textbf{estratificada}, para que las dos partes tengan la
misma proporción de clases. Recién después \textbf{ajustamos} las transformaciones: los promedios, las
columnas del one-hot y el escalador se calculan \textbf{solo con el entrenamiento}. Y eso se
\textbf{aplica} igual, sin recalcular, a validación y a test. Si el escalador o los promedios vieran la
validación, estaríamos evaluando con información que el modelo no debería tener, y las métricas
saldrían infladas.
\end{speech}

\begin{center}
\begin{tikzpicture}[caja/.style={draw=azul, fill=fondoazul, rounded corners=2pt, align=center,
    font=\small, minimum width=2.6cm, minimum height=0.9cm},
  flecha/.style={-{Stealth[length=2mm]}, thick, draw=azul}]
  \node[caja] (ent) at (0,0) {Entrenamiento\\(10.000)};
  \node[caja, right=1.4cm of ent] (aj) {Calcular promedios,\\one-hot, escalador};
  \node[caja, right=1.4cm of aj] (val) {Validación (2.500)\\y test (500)};
  \draw[flecha] (ent) -- node[above, font=\scriptsize] {fit} (aj);
  \draw[flecha] (aj) -- node[above, font=\scriptsize] {transform} (val);
\end{tikzpicture}
\end{center}

\begin{extras}[PARA ENTENDERLO]
\begin{itemize}
  \item Es el mismo patrón \texttt{fit\_transform} / \texttt{transform} del laboratorio 3. En el notebook,
        \texttt{preprocesar} y \texttt{escalar} reciben primero el conjunto con el que se ajusta y
        después los demás.\nb{3.4, 5.1}
  \item \textbf{Validación} sirve para elegir entre experimentos; \textbf{test} no se usa para decidir nada.
  \item Las transformaciones fila a fila (extraer el año, ceros a faltante) no aprenden nada de los datos,
        así que no generan fuga y se pueden hacer antes de dividir.\nb{3.1}
  \item Para el modelo final, todo se vuelve a ajustar con el train completo (12.500) y se aplica al test.
\end{itemize}
\end{extras}

% ------------------------------------------------------------------
\slide{8}{El árbol por defecto memoriza}{40 s}

\begin{speech}
Arrancamos con el \textbf{árbol de decisión con los parámetros por defecto}. Sin límites, crece hasta
profundidad 30, con 1763 hojas, y acierta el \textbf{100\,\% del entrenamiento} pero solo el
\textbf{67\,\% de validación}: memorizó los ejemplos en lugar de aprender reglas generales. Eso es
\textbf{sobreajuste}.

El gráfico muestra el F1 macro según la profundidad máxima: la curva azul, de entrenamiento, sube
siempre; la naranja, de validación, sube hasta \textbf{profundidad 10} y después se estanca y baja. Ahí
el árbol deja de aprender patrones y empieza a memorizar.
\end{speech}

\begin{extras}[PARA ENTENDERLO]
\begin{itemize}
  \item \textbf{Sobreajuste}: rinde mucho mejor en entrenamiento que en datos nuevos. \textbf{Subajuste}:
        rinde mal en los dos porque el modelo es demasiado simple (la zona izquierda del gráfico, con
        profundidad 2 a 5).
  \item Un árbol sin límites sigue dividiendo hasta que cada hoja es \textbf{pura} (una sola clase); por
        eso llega al 100\,\%. 1763 hojas para 10.000 ejemplos son menos de 6 ejemplos por hoja.
  \item Scikit-learn usa \textbf{CART}: cortes binarios ``atributo $\le$ umbral'', y elige el corte que
        más reduce la impureza de \textbf{Gini}, $1 - \sum_i p_i^2$ (criterio por defecto).\teo{10}
        ID3, el otro algoritmo de la teoría, usa la entropía y la ganancia de información.
\end{itemize}
\end{extras}

% ------------------------------------------------------------------
\slide{9}{Árbol: evolución de los experimentos}{55 s}

\begin{speech}
La tabla resume la evolución. \textbf{E1} es el árbol por defecto: F1 de 1 en entrenamiento y 0.624 en
validación. En \textbf{E2} limitamos la profundidad a 10, el mejor valor de un barrido de 2 a 20, y
validación sube a 0.665.

En \textbf{E3} agregamos \texttt{min\_samples\_split}, el mínimo de ejemplos que necesita un nodo para
poder dividirse. Con 100, el F1 de validación mejora apenas, a 0.670, pero lo importante es que
\textbf{la brecha entre entrenamiento y validación baja mucho}: el árbol rinde casi igual en datos vistos
y nuevos.

Como ajustamos un parámetro por vez, en \textbf{E4} probamos las 25 combinaciones de los dos juntos,
porque la mejor profundidad podía cambiar al exigir un mínimo. El barrido \textbf{confirmó el mínimo de
100}; profundidad 10 y 12 empataron, y elegimos 12. El árbol final tiene \textbf{162 hojas}, contra 1763
del original. Todo con semilla fija, para que sea reproducible.
\end{speech}

\begin{extras}[PARA ENTENDERLO]
\begin{itemize}
  \item Brecha de F1 entre entrenamiento y validación: E2 0.075, E3 0.020. Con
        \texttt{min\_samples\_split=200} el árbol ya subajusta (0.651 y 0.648).\nb{4.4}
  \item Grilla del barrido conjunto: profundidad 6, 8, 10, 12, 15 por mínimo 2, 10, 50, 100, 200. Con
        profundidad 6 ninguna combinación pasa de 0.583.\nb{4.5}
  \item Que dos configuraciones empaten muestra que \textbf{el resultado no depende de un valor exacto}:
        es una zona estable, no un pico casual.
  \item \texttt{min\_samples\_split} actúa como una \textbf{pre-poda}: frena el crecimiento antes, en vez
        de podar el árbol después de construido.
  \item E1 a E4 se evalúan siempre en el mismo conjunto de validación; la semilla es
        \texttt{random\_state=42}.
\end{itemize}
\end{extras}

% ------------------------------------------------------------------
\slide{10}{La popularidad decide primero}{35 s}

\begin{speech}
Una ventaja del árbol es que \textbf{se puede leer}. Estos son los primeros niveles del árbol elegido. El
primer corte es la \textbf{popularidad de la película}, y el segundo nivel también: por ejemplo, las 685
películas de entrenamiento con popularidad menor a 1.2 son \textbf{todas Malas}. Después aparece la
popularidad de las películas de los intérpretes. Esto \textbf{coincide con el análisis exploratorio}: la
popularidad era el atributo que mejor separaba las clases.
\end{speech}

\begin{extras}[PARA ENTENDERLO]
\begin{itemize}
  \item \textbf{Cómo leer cada nodo}: la condición del corte (si se cumple, se va a la izquierda);
        \texttt{samples}, cuántos ejemplos de entrenamiento llegan; \texttt{value}, cuántos hay de cada
        clase; \texttt{class}, la clase mayoritaria, que es lo que predeciría si fuera hoja.
  \item \textbf{Ojo, pregunta trampa}: el orden de \texttt{value} es \textbf{alfabético}, como
        \texttt{classes\_}: [Buena, Mala, Muy buena]. La raíz tiene [4615, 3846, 1539]: 4615 Buenas
        (46\,\%), 3846 Malas y 1539 Muy buenas.
  \item Gini de la raíz: $1 - (0{,}46^2 + 0{,}38^2 + 0{,}15^2) \approx 0{,}62$. La hoja de 685 Malas
        tiene Gini 0 (pura).
  \item Los umbrales están en la escala original de la popularidad, porque el árbol no usa escalado
        ni logaritmo.
  \item El color indica la clase mayoritaria (verde Mala, naranja Buena) y la intensidad, qué tan pura es.
  \item Se muestran solo 2 niveles (\texttt{plot\_tree(max\_depth=2)}); el árbol completo tiene 12.
\end{itemize}
\end{extras}

% ------------------------------------------------------------------
\slide{11}{Cuándo dejar de entrenar}{45 s}

\begin{speech}
Pasamos a la \textbf{red neuronal}. El modelo base son capas densas con \textbf{ReLU}, una salida
\textbf{softmax} de 3 neuronas (una probabilidad por clase), optimizador \textbf{Adam} y, como pérdida, la
\textbf{entropía cruzada categórica}, como en el laboratorio 3.

La pregunta clave es \textbf{cuántas épocas entrenar}. Para eso pasamos la validación al entrenamiento y
miramos las curvas. Este ejemplo es la red más grande, de tres capas: la pérdida de entrenamiento baja
todo el tiempo, pero la de validación \textbf{toca su mínimo en la época 4 y después sube sin parar}. Eso
es sobreajuste. El criterio fue: cada configuración se entrena primero con muchas épocas, y \textbf{la
época con menor pérdida en validación nos dice cuántas usar}.
\end{speech}

\begin{extras}[PARA ENTENDERLO]
\begin{itemize}
  \item \textbf{Época}: una pasada completa por los 10.000 ejemplos. \textbf{Batch}: cuántos ejemplos se
        procesan antes de actualizar los pesos (32 por defecto: unas 313 actualizaciones por época).
  \item \textbf{ReLU}: $\max(0, x)$. Es la activación de las capas ocultas; sin una función no lineal,
        varias capas equivaldrían a una sola transformación lineal.
  \item \textbf{Softmax}: convierte las 3 salidas en probabilidades que suman 1; la clase predicha es la
        de mayor probabilidad (\texttt{argmax}).
  \item \textbf{Entropía cruzada}: $-\log(p)$ de la probabilidad que la red le dio a la clase correcta.
        Vale 0 si la red da probabilidad 1 a la clase correcta y crece mucho cuando está segura y se
        equivoca. Necesita la clase como vector one-hot, por ejemplo Buena $\to (0, 1, 0)$.
  \item \textbf{Entrenamiento}: \textit{backpropagation} calcula cuánto contribuye cada peso al error, y
        el descenso por gradiente ajusta cada peso un paso en la dirección que lo reduce.\teo{11}
        \textbf{Adam} es una variante que adapta el tamaño del paso para cada peso.
  \item La curva de la derecha (accuracy) cuenta lo mismo: entrenamiento sube, validación baja.
\end{itemize}
\end{extras}

% ------------------------------------------------------------------
\slide{12}{Red: evolución de los experimentos}{55 s}

\begin{speech}
Esta es la evolución de la red. \textbf{R1}, una sola capa oculta de 16 neuronas con 200 épocas, llega a
0.644 de F1 en validación y \textbf{no sobreajusta}: su pérdida de validación se estabiliza y no vuelve a
subir.

Después probamos \textbf{más capacidad}. Con dos capas de 64 y 32, R2 sobreajusta desde la época 26;
cortando ahí, en R3, mantiene el F1 con mucha menos brecha, pero no supera a R1. Con tres capas, R4
sobreajusta desde la época 4, y si cortamos en 4, R5 no llega a aprender.

Partiendo de R3 variamos la \textbf{tasa de aprendizaje}, las dos filas R6: con 0.0001 todavía está
aprendiendo al final, y con 0.01 sobreajusta en la época 7. Y con \textbf{batch de 128}, R7 hace cuatro
veces menos actualizaciones y rinde algo peor. Todas quedan entre 0.594 y 0.644: \textbf{agregar
capacidad no ayudó}, y nos quedamos con la red más chica.
\end{speech}

\begin{extras}[PARA ENTENDERLO]
\begin{itemize}
  \item \textbf{Cantidad de pesos} (cada neurona tiene un peso por entrada más un sesgo):
        R1: $28 \cdot 16 + 16 + 16 \cdot 3 + 3 = 515$. R4: unos 14.000 (14.147). Con 10.000 ejemplos,
        R4 tiene más parámetros que ejemplos: por eso memoriza tan rápido.
  \item \textbf{Tasa de aprendizaje}: el tamaño de cada paso del descenso por gradiente. Muy chica, aprende
        lento; muy grande, los pasos se pasan del mínimo y el entrenamiento es inestable.
  \item \textbf{Batch más grande}: gradiente menos ruidoso pero menos actualizaciones por época
        (79 contra 313).
  \item Se variaron \textbf{uno por vez}, partiendo de la mejor configuración anterior, para saber a qué
        atribuir cada cambio.
  \item La red final es R1 \textbf{reentrenada con el train completo}, 200 épocas.\nb{5.10}
\end{itemize}
\end{extras}

% ------------------------------------------------------------------
\slide{13}{Árbol contra red, en validación}{40 s}

\begin{speech}
Comparamos los dos modelos elegidos sobre \textbf{el mismo conjunto de validación}. El árbol saca 0.723
de accuracy y \textbf{0.670 de F1 macro}; la red, 0.714 y \textbf{0.644}. En las clases frecuentes, Mala y
Buena, están \textbf{prácticamente parejos}; la red incluso es un poco mejor en Mala. La diferencia está
en \textbf{Muy buena}: F1 de 0.499 contra 0.422. Si miramos el recall, de las películas que realmente son
Muy buenas, el árbol encuentra el 39\,\% y la red el 32\,\%. Es la clase minoritaria, y la que les cuesta
a los dos.
\end{speech}

\begin{extras}[PARA ENTENDERLO]
\begin{itemize}
  \item En \textit{Muy buena} los dos tienen precisión aceptable (0.68 y 0.64) pero recall bajo: cuando
        dicen ``Muy buena'' suelen acertar, pero \textbf{la dicen pocas veces}.\nb{4.7, 5.9}
  \item Es esperable: con solo el 15\,\% de ejemplos, minimizar el error general empuja a los modelos
        hacia las clases frecuentes.
  \item La diferencia de accuracy (0.009) es chica; la de F1 macro (0.026) es más grande justamente
        porque el F1 macro le da peso a \textit{Muy buena}.
\end{itemize}
\end{extras}

% ------------------------------------------------------------------
\slide{14}{Elegimos el árbol}{35 s}

\begin{speech}
En conclusión, \textbf{elegimos el árbol}. Primero, por el número: 0.670 de F1 macro contra 0.644.
Segundo, porque es \textbf{más simple}: 162 hojas, no necesita escalado y podemos leer qué atributos usa
para decidir. Y tercero, una conclusión sobre el problema: la red más chica fue la mejor, y las grandes
sobreajustaban en pocas épocas. Eso nos indica que \textbf{el límite no está en la capacidad del modelo
sino en la información que tienen los atributos}. Gracias. ¿Preguntas?
\end{speech}

\begin{extras}[PARA ENTENDERLO]
\begin{itemize}
  \item ``El límite son los datos'' es una \textbf{interpretación}, no algo demostrado: lo sostiene que
        modelos muy distintos (árbol y red, chicos y grandes) se estancan en la misma zona de F1.
  \item El archivo entregado tiene las predicciones de los dos modelos finales sobre el test (columnas
        \texttt{pred\_arbol} y \texttt{pred\_rna}).\nb{7}
\end{itemize}
\end{extras}

% ==================================================================
\section{Preguntas posibles}

Respuestas cortas: alcanza con decir la idea y, si repreguntan, ampliar con lo de las secciones
``Para entenderlo''.

\subsection{Datos y análisis exploratorio}

\begin{description}[style=nextline, leftmargin=0pt, itemsep=6pt]
  \item[¿Por qué F1 macro y no accuracy?]
    Las clases están desbalanceadas. Un modelo que casi nunca predice \textit{Muy buena} puede tener buena
    accuracy; el F1 macro promedia las tres clases por igual y lo castiga. Igual reportamos las dos.
  \item[¿Cómo detectaron que los ceros eran datos faltantes?]
    Por sentido común del dominio (una película no cuesta 0 ni dura 0 minutos) y por la proporción: un 79\,\%
    de presupuestos en 0 no puede ser real. Los contamos aparte porque \texttt{isna()} no los ve.
  \item[¿Por qué dejaron el año si casi no separa las clases?]
    Por sí solo separa poco, pero el árbol puede combinarlo con otros atributos (por ejemplo, la popularidad
    significa cosas distintas en una película vieja y en una nueva). Y no agrega riesgo de fuga.
  \item[¿El test tiene la misma distribución de clases que el train?]
    Lo verificamos en el notebook (sección 2.8), junto con que no hay películas repetidas entre train y test.
\end{description}

\subsection{Preprocesamiento}

\begin{description}[style=nextline, leftmargin=0pt, itemsep=6pt]
  \item[¿Qué es la fuga de datos y cómo la evitaron?]
    Usar, al entrenar o transformar, información de los datos con los que después se evalúa. La evitamos
    dividiendo primero y calculando promedios, columnas del one-hot y escalador solo con entrenamiento;
    después se aplican igual a validación y test.
  \item[¿Por qué completar con el promedio y no con otro valor?]
    Es lo que vimos en el laboratorio 1 con \texttt{Age}. La mediana sería más robusta con colas largas, y
    es una mejora posible. Además, los indicadores \texttt{*\_conocido} le permiten al modelo distinguir un
    valor imputado de uno real.
  \item[Si el 79\,\% del presupuesto falta, ¿no es inventar datos completarlo?]
    Sí, por eso agregamos \texttt{budget\_conocido}: el árbol puede separar primero si el dato existe y
    recién después usar el valor. Con el país no lo hicimos porque es categórico y falta todavía más
    (90\,\%): completarlo con el más frecuente sería inventar casi todo.
  \item[¿Por qué escalar para la red y no para el árbol?]
    El árbol compara cada atributo contra un umbral: cambiar la escala solo cambia el número del umbral,
    no qué ejemplos quedan de cada lado. La red multiplica las entradas por pesos y aprende por gradiente:
    con escalas muy distintas (millones contra 0/1) el entrenamiento es lento e inestable.
  \item[¿Por qué el logaritmo antes de escalar?]
    \texttt{MinMaxScaler} usa el mínimo y el máximo. Con colas largas, el máximo lo definen unos pocos
    valores enormes y el resto queda aplastado cerca de $-1$. Sin el logaritmo, la misma red bajó de 0.65 a
    0.44 de F1 macro.
  \item[{¿Por qué al intervalo $[-1, 1]$?}]
    Es como lo hicimos en el laboratorio 3. Lo importante es que todas las entradas queden en un rango
    comparable; \texttt{StandardScaler} (media 0, desvío 1) también hubiera servido.
  \item[¿Por qué descartaron los nombres de actores y directores?]
    Miles de valores distintos: el one-hot crearía miles de columnas casi vacías y el modelo memorizaría
    nombres que en el test casi no se repiten. Su popularidad y cantidad de películas ya están como números.
  \item[¿Qué pasa si en el test aparece un género que no estaba en el entrenamiento?]
    Las columnas del one-hot se fijan con el conjunto de ajuste (\texttt{reindex}): un género nuevo queda
    con todas sus columnas en 0, y una columna que falte se completa con ceros.
  \item[¿Por qué estratificar la división?]
    Para que entrenamiento y validación tengan la misma proporción de clases. Con \textit{Muy buena} al
    15\,\%, una división al azar podría dejar a validación con más o menos de esa clase y sesgar las
    métricas.
\end{description}

\subsection{Árbol de decisión}

\begin{description}[style=nextline, leftmargin=0pt, itemsep=6pt]
  \item[¿Cómo elige el árbol el atributo y el umbral de cada corte?]
    Prueba los cortes posibles ``atributo $\le$ umbral'' y elige el que deja los hijos más puros, es decir,
    el que más reduce la impureza de Gini (criterio por defecto de scikit-learn, algoritmo CART). ID3 hace
    lo mismo con la entropía (ganancia de información).
  \item[¿Qué hace \texttt{max\_depth}? ¿Y \texttt{min\_samples\_split}?]
    \texttt{max\_depth} limita cuántas preguntas seguidas puede hacer el árbol. \texttt{min\_samples\_split}
    exige que un nodo tenga al menos esa cantidad de ejemplos para dividirse: evita reglas apoyadas en
    muy pocos casos. Los dos frenan el sobreajuste.
  \item[¿Por qué un parámetro por vez y después el barrido conjunto?]
    Un parámetro por vez permite ver el efecto de cada uno. Pero la mejor profundidad podía cambiar al
    exigir un mínimo para dividir; el barrido conjunto lo verificó y confirmó la elección.
  \item[¿Cómo predice una hoja?]
    Con la clase mayoritaria de los ejemplos de entrenamiento que llegaron a esa hoja.
  \item[En la figura del árbol, ¿qué es \texttt{value = [4615, 3846, 1539]}?]
    Cuántos ejemplos de cada clase llegan al nodo, en orden \textbf{alfabético}: Buena, Mala, Muy buena.
  \item[¿Por qué el árbol por defecto llega al 100\,\% en entrenamiento?]
    Porque sigue dividiendo hasta que cada hoja tiene una sola clase: memoriza cada ejemplo.
  \item[¿Por qué eligieron profundidad 12 si empata con 10?]
    Fue la de mayor F1 en validación del barrido: con más decimales, 12 queda apenas por encima de 10, y
    redondeadas las dos dan 0.670. El empate muestra que el resultado es estable; con 10 el árbol sería
    algo más chico y también sería una elección válida.
\end{description}

\subsection{Red neuronal}

\begin{description}[style=nextline, leftmargin=0pt, itemsep=6pt]
  \item[¿Por qué softmax en la salida y \texttt{categorical\_crossentropy}?]
    Softmax da 3 probabilidades que suman 1, una por clase. La entropía cruzada mide qué tan baja fue la
    probabilidad asignada a la clase correcta, y necesita la clase como vector one-hot.
  \item[¿Para qué sirve ReLU?]
    Es la no linealidad de las capas ocultas. Sin ella, apilar capas sería equivalente a una sola
    transformación lineal.
  \item[¿Cómo aprende la red?]
    Hace la pasada hacia adelante, calcula la pérdida, y con backpropagation obtiene el gradiente de la
    pérdida respecto de cada peso; el optimizador (Adam) mueve los pesos un paso en la dirección que
    reduce la pérdida. Eso se repite en cada batch.
  \item[¿Cómo eligieron las épocas?]
    Mirando la curva de pérdida en validación: cuando sube mientras la de entrenamiento sigue bajando, la
    red sobreajusta. Usamos la época con menor pérdida en validación.
  \item[¿No es eso hacer trampa con validación?]
    Validación existe justamente para elegir hiperparámetros, y las épocas son uno. Lo que no se toca para
    decidir es el test.
  \item[¿Por qué la red más chica fue la mejor?]
    Con 28 atributos y 10.000 ejemplos, las redes grandes tienen muchísimos pesos (R4 más de 14.000) y
    memorizan en pocas épocas. Cortadas a tiempo, no aprenden más que la chica: la información disponible
    en los atributos es la que limita.
  \item[¿Qué es una época y qué es el batch?]
    Época: una pasada por todo el entrenamiento. Batch: cuántos ejemplos se procesan antes de cada
    actualización de pesos.
  \item[¿Los resultados de la red son reproducibles?]
    Sí: fijamos la semilla (\texttt{keras.utils.set\_random\_seed}) y activamos el modo determinista de
    TensorFlow, así cada corrida da los mismos pesos iniciales y el mismo orden de batches.
  \item[¿Por qué no usaron Dropout o early stopping?]
    Nos quedamos con lo visto en la teoría y los laboratorios. El efecto del early stopping lo
    reproducimos a mano eligiendo las épocas con la curva de validación.
\end{description}

\subsection{Resultados y cierre}

\begin{description}[style=nextline, leftmargin=0pt, itemsep=6pt]
  \item[¿Cómo les fue en el test?]
    El test trae la columna \texttt{class}; la usamos \textbf{una sola vez al final}, solo para informar,
    sin cambiar nada. Árbol: accuracy 0.704 y F1 macro 0.669. Red: 0.696 y 0.622. Confirma el orden de
    validación.\nb{7.1}
  \item[¿Por qué a la red le va peor en \textit{Muy buena}?]
    Predice esa clase muy pocas veces: en el test, para el 6\,\% de las películas, cuando en realidad son
    el 15\,\%. Su recall es bajo.
  \item[¿Qué harían para mejorar \textit{Muy buena}?]
    Darle más peso a esa clase al entrenar (por ejemplo, \texttt{class\_weight}), y mejorar los atributos:
    es donde vemos el límite.
  \item[¿Qué más probarían?]
    Validación cruzada para tener métricas más estables que con una sola división; varios árboles
    combinados (la teoría menciona técnicas que construyen más de un árbol, como random forest);
    completar faltantes con la mediana.
  \item[¿Por qué el modelo final se reentrena con todo el train?]
    Validación ya cumplió su función (elegir hiperparámetros); con la configuración fija, entrenar con
    12.500 ejemplos en vez de 10.000 suele dar un modelo algo mejor.
  \item[¿Cómo usaron la IA en el trabajo?]
    Como está en la sección 8 del notebook: para consultar cómo llevar a pandas, scikit-learn y Keras lo
    visto en clase y en los laboratorios, para los gráficos y para el diseño de las slides. Las decisiones y
    los experimentos los podemos justificar uno por uno.
\end{description}

\section{Puntos débiles y cómo responderlos}

Mejor conocerlos antes de que los señalen. En todos, la respuesta honesta es la mejor.

\begin{aviso}{Lo que pueden señalar}
\begin{itemize}
  \item \textbf{Hay dos filas ``R6'' en la tabla de la red.} Son las dos variantes de tasa de aprendizaje
        del mismo experimento (0.0001 y 0.01).
  \item \textbf{La slide 11 dice ``la época 4 de 100'' y la red elegida usa 200 épocas.} La curva es de R4
        (la red grande), como ejemplo de sobreajuste. R1 no sobreajusta: su pérdida de validación se
        estabiliza, así que 200 épocas no la perjudican.
  \item \textbf{Las métricas de validación son algo optimistas.} Elegimos el mejor de varios experimentos
        con el mismo conjunto de validación, y eso favorece al elegido. Por eso el test, que no se usó para
        decidir, da algo menos (0.669 contra 0.670 en el árbol; 0.622 contra 0.644 en la red).
  \item \textbf{Las diferencias son chicas.} Entre las configuraciones de la red hay pocos centésimos, y
        con una sola división no sabemos cuánto de eso es ruido. La ventaja del árbol sobre la red, en
        cambio, se repite en validación y en test.
  \item \textbf{El promedio con colas largas.} El promedio del presupuesto lo tiran hacia arriba unos pocos
        valores enormes; la mediana sería más representativa. Lo mitiga el indicador \texttt{*\_conocido}.
  \item \textbf{¿La popularidad es una ``trampa''?} Es un atributo dado en los datos, no la clase. Pero
        puede estar medida después del estreno, y entonces reflejar en parte lo mismo que la calificación.
        Es una limitación del dataset que vale la pena reconocer si preguntan.
  \item \textbf{\textit{Muy buena} sigue siendo débil} (recall de 39\,\% y 32\,\%). No aplicamos ninguna
        técnica de balanceo; es la primera mejora que haríamos.
\end{itemize}
\end{aviso}

\section{Números para tener a mano}

\begin{center}
\small
\begin{tabularx}{\linewidth}{@{}l X@{}}
\toprule
Dato & Valor \\
\midrule
Train / test                    & 12.500 / 500 películas \\
División                        & 10.000 entrenamiento / 2.500 validación, estratificada, semilla 42 \\
Clases                          & Mala 38.5\,\%, Buena 46.2\,\%, Muy buena 15.4\,\% \\
Ceros y faltantes               & budget 79\,\%, revenue 76\,\%, país 90\,\%, runtime 10\,\% \\
Atributos de entrada            & 28 \\
Árbol por defecto               & profundidad 30, 1763 hojas; accuracy 100\,\% / 67\,\% \\
Árbol elegido                   & \texttt{max\_depth=12}, \texttt{min\_samples\_split=100}, 162 hojas \\
Red elegida                     & 1 capa oculta de 16 (ReLU), softmax 3, Adam lr 0.001, batch 32, 200 épocas, 515 pesos \\
Validación (F1 macro)           & árbol 0.670, red 0.644 \\
Validación (accuracy)           & árbol 0.723, red 0.714 \\
F1 \textit{Muy buena}           & árbol 0.499, red 0.422 (recall 39\,\% y 32\,\%) \\
Test (F1 macro, solo informativo) & árbol 0.669, red 0.622 \\
\bottomrule
\end{tabularx}
\end{center}

\end{document}
